\chapter{Column- and Row-Level Security, Data Tagging, and Erasure}
\index{masking policy}\index{row access policy}\index{object tagging}\index{GDPR}\index{Week 18}%
\begin{objectives}
\begin{itemize}
    \item Classify sensitive data with object tags on databases, schemas, tables, and columns.
    \item Enforce dynamic data masking policies keyed on session context and tags.
    \item Restrict row visibility with row access policies driven by \texttt{CURRENT\_ROLE()} and entitlements tables.
    \item Propagate GDPR erasure across schemas, clones, streams, and shares---or erase by crypto-shredding.
    \item Build a hands-on tag-based masking policy plus a regional row access policy.
    \item Design a HIPAA/GDPR compliance framework with automatic classification and dynamic masking.
\end{itemize}
\end{objectives}

\begin{crosscloud}
Classification-driven policy enforcement is the industry pattern: tag the data once, let policies follow the tags everywhere.

\vspace{2mm}
{\footnotesize
\begin{tabular}{@{}p{2.8cm}p{3.0cm}p{3.0cm}p{3.0cm}@{}}
\toprule
\textbf{Capability} & \textbf{AWS} & \textbf{Azure / Fabric} & \textbf{GCP (BigQuery)} \\
\midrule
Classification & Macie + resource tags & Purview sensitivity labels & Sensitive Data Protection + policy tags \\
Column masking & Redshift dynamic masking / Lake Formation & Synapse/Fabric dynamic masking & BigQuery dynamic masking via policy tags \\
Row-level security & Lake Formation row filters & Synapse RLS / OneLake security & BigQuery row access policies \\
Policy-tag binding & LF tag expressions & Label-inherited protection & Policy tag taxonomy \\
Erasure support & S3 Object Lock exceptions, KMS shred & Purview + CMK shredding & CMEK key destruction \\
\addlinespace
Snowflake equivalent & \multicolumn{3}{l}{Object tags, tag-based masking policies, row access policies, crypto-shredding} \\
\bottomrule
\end{tabular}}

\vspace{1mm}
\textbf{MongoDB} offers client-side field-level encryption; \textbf{Fabric} inherits Purview labels into OneLake. Snowflake's distinctive leverage: tags attach to \emph{columns}, policies attach to \emph{tags}, so classifying a column once protects it in every query, clone, and share that touches it.
\end{crosscloud}

\section{Week 18 Roadmap}
Chapter 17 controlled who reaches the door. This chapter controls what they see through it: masking columns for unauthorized eyes, filtering rows by entitlement, and---when regulation demands---making data vanish everywhere it ever flowed, including clones and Fail-safe \cite{snowflakeMaskingDocs,snowflakeRowAccessDocs}.

Monday covers object tagging as classification infrastructure. Tuesday builds dynamic masking. Wednesday handles row access policies and the GDPR erasure problem. Thursday implements tag-based masking plus regional row filtering. Friday assembles the full compliance framework.

\begin{keyterms}
\textbf{Object tag}: a key-value label applied to Snowflake objects for classification and governance.\quad
\textbf{Masking policy}: a column-attached rule returning obfuscated values based on session context.\quad
\textbf{Tag-based masking}: a masking policy attached to a tag, protecting every column carrying that tag.\quad
\textbf{Row access policy}: a table-attached predicate filtering visible rows by caller context.\quad
\textbf{Crypto-shredding}: rendering data unrecoverable by destroying its per-subject encryption key.
\end{keyterms}

\section{Monday: Object Tagging as Classification Infrastructure}
\index{object tagging}\index{data classification}
An object tag is governed metadata: \texttt{pii\_email = 'true'} on a column declares what the data \emph{is}, independent of who queries it. Tags are schema objects with their own RBAC, propagate to columns via \texttt{ALTER TABLE ... SET TAG}, and are visible in \texttt{ACCOUNT\_USAGE.TAG\_REFERENCES}---which turns ``where is all the PII?'' from an audit project into a query.

\begin{definitionbox}
\textbf{Tag-based governance} separates classification from enforcement: tags classify data once; masking policies, access rules, and lineage reports all consume the tags. Adding a new PII column is a tagging decision, and every existing control applies automatically.
\end{definitionbox}

Snowflake's \texttt{SNOWFLAKE.CORE} functions (and automatic classification for new tables) can suggest PII tags by profiling column content---useful for the initial estate inventory, but the engineered state is intentional tags in DDL, reviewed like code. Tags also power lineage: \texttt{OBJECT\_TAGGING} lineage reports show where tagged data flows through views, CTAS, and clones.

\begin{tipbox}
Standardize a tag taxonomy early: \texttt{pii\_*} (classification), \texttt{fin\_*} (financial), \texttt{retention\_*} (lifecycle). Allowed-values lists on tags (\texttt{ALLOWED\_VALUES}) prevent the taxonomy from decaying into free text.
\end{tipbox}

\section{Tuesday: Dynamic Data Masking}
\index{masking policy}\index{tag-based masking}
A masking policy is a function over a column value plus the caller's context: return the value for entitled roles, an obfuscation for everyone else. Policies evaluate per query at runtime---the stored data is untouched, so no desensitized copies drift out of sync.

\begin{lstlisting}[language=SQL, caption={A role-aware masking policy, bound to a tag}]
CREATE OR REPLACE MASKING POLICY pii_mask AS
  (val STRING) RETURNS STRING ->
  CASE WHEN IS_ROLE_IN_SESSION('HR_ROLE') THEN val
       WHEN IS_ROLE_IN_SESSION('ANALYST') THEN
            REGEXP_REPLACE(val, '.+@', '***@')   -- domain visible
       ELSE '***MASKED***' END;

ALTER TAG pii_email SET MASKING POLICY pii_mask;
-- every column tagged pii_email is now masked by role
\end{lstlisting}

Binding the policy to the \emph{tag} rather than each column is the scaling move: the hundredth PII column is protected by tagging it, with zero new policy code. Masking policies apply to views, clones, and shares automatically---the policy follows the column, not the container.

\begin{pitfall}
\textbf{Masking does not protect against aggregate inference.} If analysts can \texttt{AVG(salary) GROUP BY department} on unmasked numerics, small groups leak individuals. Masking is column-level; inference control needs row access policies (Wednesday) plus minimum-cohort rules in shared products (Chapter 20).
\end{pitfall}

\begin{notebox}
Test masking policies against \emph{every consuming role}, including simulated \texttt{ACCOUNTADMIN}. Policy bugs that show cleartext to an over-privileged role are found by negative tests, not by reading the CASE expression.
\end{notebox}

\section{Wednesday: Row Access Policies and GDPR Erasure}
\index{row access policy}\index{GDPR erasure}\index{crypto-shredding}
\subsection{Row Access Policies}
A row access policy is a filter the engine applies invisibly: the query's author writes \texttt{SELECT * FROM sales}; the engine evaluates the policy per row. Entitlements typically live in a mapping table---\texttt{sales\_entitlements(role, region)}---so policy bodies stay stable while access changes are data changes:

\begin{lstlisting}[language=SQL, caption={Row access policy driven by an entitlements table}]
CREATE OR REPLACE ROW ACCESS POLICY regional_rap AS
  (sales_region STRING) RETURNS BOOLEAN ->
  IS_ROLE_IN_SESSION('REGIONAL_ADMIN')
  OR EXISTS (
    SELECT 1 FROM gov.sales_entitlements e
    WHERE e.role_name = CURRENT_ROLE()
      AND e.region = sales_region
  );

ALTER TABLE gold.sales ADD ROW ACCESS POLICY regional_rap ON (region);
\end{lstlisting}

Policies compose with masking orthogonally: masking decides \emph{what a visible row shows}; row policies decide \emph{which rows exist for you}. Multi-tenant designs lean on both.

\subsection{GDPR Erasure: Why DELETE Is Not Enough}
\index{right to erasure}
A naive \texttt{DELETE} leaves the subject recoverable everywhere: Time Travel (up to 90 days), Fail-safe (7 further days), clones created before the delete (they share micro-partitions), unconsumed stream offsets (deletes replay downstream), Snowpipe Streaming channel buffers, and reader-account views. A defensible erasure process therefore chooses one of two architectures:

\begin{enumerate}
    \item \textbf{Retention-bounded deletion:} delete the rows everywhere (raw, silver, gold; recreate or swap clones; let stream consumers replay the delete), set \texttt{DATA\_RETENTION\_TIME\_IN\_DAYS = 0} on PII schemas, and wait out Fail-safe---documented, slow, complete.
    \item \textbf{Crypto-shredding:} encrypt PII per data subject with subject-specific keys held in an external KMS (reached via external functions, or customer-managed keys with Tri-Secret Secure). Erasure destroys the subject's key; every copy---clones, backups, Fail-safe---becomes permanently unreadable instantly.
\end{enumerate}

\begin{pitfall}
\textbf{Legal hold must override erasure.} An erasure task checks a legal-hold table first; held subjects are exempted with extended \texttt{DATA\_RETENTION\_TIME\_IN\_DAYS} documented as the audit-trail exception. Deleting held data is worse than keeping it.
\end{pitfall}

\section{Thursday Hands-On Project: Tag-Based Masking and Regional Row Filtering}
\index{hands-on project!masking}
This lab implements both enforcement layers and tests them adversarially.

\subsection*{Scenario}
HR data carries SSNs that only HR may see; sales reps must see only their region's rows. New PII columns added next quarter must be protected by tagging alone.

\subsection*{Procedure}
\begin{enumerate}
    \item Create tags \texttt{pii\_ssn} and \texttt{pii\_email} with the governance taxonomy.
    \item Create the masking policy from Tuesday and bind it to both tags; tag the relevant columns.
    \item Create the entitlements table and regional row access policy from Wednesday; attach to \texttt{gold.sales}.
    \item As each role (HR, ANALYST, SALES\_EAST, SALES\_WEST), run the same queries and record the visible differences.
    \item Add a brand-new SSN column, tag it only, and prove masking applies with no policy change.
    \item Attempt bypasses: aggregate inference on a small group, and a direct entitlements-table edit attempt as a non-admin.
\end{enumerate}

\subsection*{Deliverables}
\begin{itemize}
    \item \textbf{DDL pack:} tags, masking policy, row access policy, entitlements table, grants.
    \item \textbf{Test matrix:} role $\times$ query with observed outputs, including the bypass attempts.
    \item \textbf{Classification report:} \texttt{TAG\_REFERENCES} query listing all PII columns and their policies.
    \item \textbf{Rollout plan:} how new columns acquire protection (tagging in CI) and how policies are recertified.
\end{itemize}

\subsection*{Acceptance Criteria}
\begin{itemize}
    \item Non-HR roles never see cleartext SSN through any tested path, including via views and clones.
    \item Regional filtering holds for direct queries, joins, and aggregates.
    \item A newly tagged column is masked with zero policy edits.
    \item The aggregate-inference test is documented with its row-policy mitigation.
\end{itemize}

\subsection*{Stretch Goals}
\begin{itemize}
    \item Implement the retention-bounded erasure procedure end to end on a sandbox subject, including clone swap.
    \item Prototype crypto-shredding with per-subject keys via an external KMS mock.
    \item Build the masking-policy CI gate: untagged PII-looking column names fail the deploy.
\end{itemize}

\section{Friday Real-World Architecture: HIPAA/GDPR Compliance Framework}
\index{Real-World Architecture!compliance framework}\index{HIPAA}
A health-tech company processes patient-adjacent data across 40 schemas. Regulators require demonstrable classification, enforced minimization, and provable erasure.

\begin{figure}[H]
    \centering
    \resizebox{\textwidth}{!}{%
    \begin{tikzpicture}[
        cls/.style={rectangle, rounded corners, draw=primary, thick, fill=primary!6, align=center, minimum width=3.0cm, minimum height=1.0cm, font=\small\bfseries},
        pol/.style={rectangle, rounded corners, draw=codegreen!60!black, thick, fill=green!8, align=center, minimum width=3.0cm, minimum height=1.0cm, font=\small\bfseries},
        era/.style={rectangle, rounded corners, draw=red!60!black, thick, fill=red!8, align=center, minimum width=3.0cm, minimum height=1.0cm, font=\small\bfseries},
        arrow/.style={-{Stealth[length=2.5mm]}, draw=primary, thick},
        lbl/.style={font=\scriptsize\itshape, text=codegray}
    ]
        \node[cls] (tag) at (0, 0) {Tag taxonomy\\pii\_*, phi\_*, retention\_*};
        \node[pol] (mask) at (4.6, 1.3) {Tag-based masking\\policies};
        \node[pol] (rap) at (4.6, -1.3) {Row access policies\\+ entitlements};
        \node[era] (shred) at (9.4, 0) {Erasure service\\crypto-shredding\\+ legal-hold check};
        \node[cls] (audit) at (13.4, 0) {Audit evidence\\TAG\_REFERENCES,\\ACCESS\_HISTORY};
        \draw[arrow] (tag) -- (mask);
        \draw[arrow] (tag) -- (rap);
        \draw[arrow] (tag) -- node[lbl, above]{subject keys} (shred);
        \draw[arrow] (mask) -- (audit);
        \draw[arrow] (rap) -- (audit);
        \draw[arrow] (shred) -- (audit);
    \end{tikzpicture}%
    }
    \caption{Compliance framework: classification tags feed enforcement (masking, row policies) and erasure (crypto-shredding), with ACCOUNT\_USAGE as the evidence layer.}
    \label{fig:chapter18-compliance}
\end{figure}

\subsection{Framework Components}
\begin{enumerate}
    \item \textbf{Classification as code:} tags in DDL, suggested tags from automatic classification reviewed into the taxonomy, and the CI gate rejecting untagged sensitive-looking columns.
    \item \textbf{Enforcement:} tag-based masking for identifiers and contact data; row access policies for ward/clinic tenancy; minimum-cohort rules on any shared aggregate.
    \item \textbf{Erasure:} per-subject envelope encryption keys in an external KMS; the erasure service checks legal hold, then destroys the key---rendering every copy unreadable, Fail-safe included.
    \item \textbf{Evidence:} \texttt{TAG\_REFERENCES} for classification coverage, \texttt{ACCESS\_HISTORY} for who touched PHI, masking-policy test results, and erasure completion logs---the audit binder regenerates itself weekly.
\end{enumerate}

\begin{pitfall}
\textbf{A framework measured by policy count instead of test results is compliance theater.} The artifacts that survive audits are the negative-test matrix, the erasure completion log, and the coverage query showing zero untagged sensitive columns---all produced continuously, not assembled before the audit.
\end{pitfall}

\section{Interview Q\&A}
\begin{enumerate}
    \item \textbf{Q: What session context does a masking policy evaluate?}\\
    A: Functions like \texttt{CURRENT\_ROLE()} and \texttt{IS\_ROLE\_IN\_SESSION()} over the querying session, evaluated per query at runtime against the stored value.
    \item \textbf{Q: How does tag-based masking scale across hundreds of columns?}\\
    A: The policy binds to the tag once; every column carrying the tag is protected automatically, including columns added later.
    \item \textbf{Q: What do row access policies typically filter on?}\\
    A: Caller context (role, user) joined to an entitlements mapping table, so access changes are data changes, not policy redeploys.
    \item \textbf{Q: Why doesn't a plain DELETE satisfy a GDPR erasure request?}\\
    A: Deleted rows persist in Time Travel, Fail-safe, pre-delete clones, unconsumed stream offsets, and shared/reader-account views.
    \item \textbf{Q: How does crypto-shredding erase data held in clones and Fail-safe?}\\
    A: PII is encrypted per subject with external KMS keys; destroying the subject's key makes every copy of the ciphertext---backups and Fail-safe included---permanently unreadable.
\end{enumerate}

\section{Further Reading}
Snowflake's column-level and row-level security guides define masking and row access policies \cite{snowflakeMaskingDocs,snowflakeRowAccessDocs}; the object-tagging documentation covers tags and lineage \cite{snowflakeTaggingDocs}. Chapter 4's Time Travel and Fail-safe mechanics define the erasure boundary \cite{snowflakeTimeTravelDocs,snowflakeFailSafeDocs}, and Chapter 20 extends these ideas into multi-party clean rooms.

\section*{Key Takeaways}
\begin{itemize}
    \item Tags classify once; policies enforce everywhere the data flows---columns, views, clones, shares.
    \item Masking decides what a visible row shows; row access policies decide which rows exist.
    \item Entitlements tables keep policy bodies stable while access changes as data.
    \item Masking is not inference control---aggregates leak; add row policies and cohort minimums.
    \item DELETE is not erasure: Time Travel, Fail-safe, clones, streams, and shares all retain.
    \item Crypto-shredding with per-subject keys is the only instant, total erasure; legal hold overrides everything.
\end{itemize}

\section{Exercises}
\begin{enumerate}
    \item \textbf{(Conceptual)} Compare masking policies, secure views, and row access policies as enforcement points.
    \item \textbf{(Conceptual)} Explain why policy-on-tag beats policy-on-column at 500 columns.
    \item \textbf{(Hands-on)} Build a masking policy with three role tiers and prove each tier's output.
    \item \textbf{(Hands-on)} Demonstrate the aggregate-inference leak and close it with a row policy or cohort floor.
    \item \textbf{(Hands-on)} Execute the retention-bounded erasure on a test subject and document every copy eliminated.
    \item \textbf{(Design)} Write the tag taxonomy and allowed-values design for a financial-services tenant.
    \item \textbf{(Design)} Design the legal-hold table and the erasure task's decision flow.
    \item \textbf{(Operations)} Build the classification-coverage report: sensitive-looking columns lacking tags.
    \item \textbf{(Security)} Threat-model the entitlements table: who can write it, and how are edits audited?
    \item \textbf{(Interview)} Explain crypto-shredding to a regulator in five sentences.
\end{enumerate}
